% STATUS: DRAFT
\chapter{Connes classification and where type III comes from}\label{ch:connes_class}

\begin{physmotiv}
Types I and II are familiar: matrices, and infinite-temperature limits of matrices. Type III looks exotic until one sees where it comes from: an \emph{infinite} system in a state that is \emph{not} maximally mixed on any finite part. Each site contributes a modular spectrum $\{\lambda,1,\lambda^{-1}\}$ and the infinite tensor product has a modular operator whose spectrum is the multiplicative group generated by all such ratios. Connes' classification uses precisely that spectrum. The refined type III$_\lambda$ ($0\le\lambda\le1$) then tells how rich the modular flow is, and in the end only the largest, III$_1$, with continuous modular spectrum, is relevant for QFT. At the end of the chapter, the Connes--Takesaki theorem tells how III$_1$ is built from the type II$_\infty$ algebras on which the rest of the book depends.
\end{physmotiv}

\section{Infinite product construction}

Consider $\bigotimes_{j=1}^\infty(M_{k_j},\varphi_j)$ with $\varphi_j=\Tr(\rho_j\,\cdot)$ and GNS vector $\Omega=\bigotimes\rho_j^{1/2}$ in $\HH=\bigotimes_j(\C^{k_j}\otimes\overline{\C^{k_j}})$. The weak closure $\M$ of the local algebra is an ``infinite tensor product of finite factors'' (ITPFI) and $\Omega$ is cyclic and separating. By \cref{ch:tomita} (example 4) the modular operator is the infinite product
\begin{equation}
\Delta=\bigotimes_j(\rho_j\otimes\rho_j^{-1}),
\end{equation}
with the local spectra $\{p_a/p_b\}$ ratios of eigenvalues of $\rho_j$. Examples:

\begin{center}
\renewcommand{\arraystretch}{1.3}
\begin{tabular}{@{}lll@{}}
\toprule
State & Spectrum of $\Delta$ (non-zero part) & Type\\
\midrule
$\rho_j=\id/k$ (trace) & $\{1\}$ & II$_1$\\
$\rho_j=\mathrm{diag}(p,1-p)$, $\lambda=\tfrac{p}{1-p}\ne1$ & $\{\lambda^n:n\in\Z\}$ & III$_\lambda$\\
two incommensurate ratios, $\lambda_1^m\lambda_2^n$, $\log\lambda_1/\log\lambda_2\notin\mathbb Q$ & dense in $\R_+$ & III$_1$ (Araki--Woods)\\
\bottomrule
\end{tabular}
\end{center}
(It is important that the infinite product is of non-trivial factors infinitely often; a finite number of ratios would give the point spectrum of a type I algebra.) The ITPFI factors $R_\lambda$ (Powers, $0<\lambda<1$) and $R_\infty$ (Araki--Woods) are the standard model examples of type III.

\section{The Connes spectrum and the classification}

Let $\M$ be a type III factor. By \cref{ch:modham} the modular flows of any two faithful normal states are inner-equivalent, so the following invariants are independent of the state.

\begin{definition}
The \emph{Connes invariant} $S(\M)\subset[0,\infty)$ is the intersection over all faithful normal states (weights) $\varphi$ of the spectrum $\spec(\Delta_\varphi)$. The \emph{$T$-invariant} $T(\M)\subset\R$ is the set of $t$ for which $\sigma^\varphi_t$ is inner.
\end{definition}

\begin{theorem}[Connes]
For a type III factor, $S(\M)\setminus\{0\}$ is a closed subgroup of $\R_{>0}$ (under multiplication) and (up to $0$) takes one of the forms
\begin{center}
\renewcommand{\arraystretch}{1.3}
\begin{tabular}{@{}lll@{}}
\toprule
Type & $S(\M)$ & $T(\M)$\\
\midrule
III$_0$ & $\{0,1\}$ & no simple formula (classified by an ergodic flow, the flow of weights)\\
III$_\lambda$, $0<\lambda<1$ & $\{0\}\cup\{\lambda^n:n\in\Z\}$ & $\frac{2\pi}{|\log\lambda|}\Z$\\
III$_1$ & $[0,\infty)$ & $\{0\}$\\
\bottomrule
\end{tabular}
\end{center}
\end{theorem}
(The latter two are \emph{definitions} of the subtypes. Type III$_\lambda$ for $0<\lambda<1$ means the modular flow is periodic modulo inner automorphisms.) For separable, \emph{hyperfinite} factors the classification is complete (Connes for $\lambda\ne1$; Haagerup 1987 for III$_1$): there is exactly one hyperfinite factor of each type III$_\lambda$, $0<\lambda\le1$.

\begin{whatwrong}
Do not read ``III$_1$'' as ``rank one'': the subscript is the ratio $\lambda$. III$_1$ means that all positive real numbers appear in the modular spectrum, equivalently the modular flow is \emph{aperiodic}, with no choice of states making $\sigma_t$ periodic modulo inner. Rindler wedges are of this kind: the boost spectrum is all of $\R$, so $\Delta=\ee^{-2\pi B}$ has spectrum $[0,\infty)$.
\end{whatwrong}

\section{The thermodynamic origin: why III appears with finite temperature}

Physically, the type III$_\lambda$ ITPFI factor arises from a non-interacting spin chain at \emph{finite} temperature with level spacing $h$: $\rho_j=\mathrm{diag}(\ee^{-\beta h/2},\ee^{\beta h/2})/(2\cosh\tfrac{\beta h}2)$, so $\lambda=\ee^{-\beta h}$ and the modular period is $2\pi/(\beta h)$. In finite volume $N$ the Gibbs state is $\Tr(\rho_N\,\cdot)$ with $\rho_N=\ee^{-\beta H_N}/Z_N$; its modular flow $\rho_N^{it}\,\cdot\,\rho_N^{-it}$ is \emph{inner}, generated by $\rho_N\in M_{2^N}$, i.e.\ by the Hamiltonian. As $N\to\infty$ the eigenvalue ratios $\ee^{-\beta hn}$ ($n\in\Z$ counting spin flips) survive and fill the group $\lambda^\Z$, but the would-be generator, the infinite-volume Hamiltonian, is no longer an element of the algebra (\cref{ch:quasilocal}); the flow becomes outer and the trace is lost. A QFT wedge is the same phenomenon with a continuum of degrees of freedom: the modular Hamiltonian $2\pi B$ has spectrum all of $\R$, so $\Delta$ has spectrum $[0,\infty)$: type III$_1$ (\cref{ch:typeIII_local}).

\section{Connes--Takesaki: III$_1$ from II$_\infty$}

The key structural result for our purposes was obtained by Connes and Takesaki.

\begin{theorem}[Connes--Takesaki]
Let $\M$ be a $\sigma$-finite von Neumann algebra of type III and $\sigma^\varphi$ the modular flow of a faithful normal state $\varphi$. Then the crossed product $\N=\M\rtimes_{\sigma^\varphi}\R$ (\cref{ch:crossed}) is a von Neumann algebra of type II$_\infty$. It carries a canonical trace $\tau$ and a \emph{dual action} $\theta_s$ of $\R$ which rescales the trace, $\tau\circ\theta_s=\ee^{-s}\tau$. The original algebra is recovered as the fixed points of the dual action, $\M=\N^{\theta}$, and (Takesaki duality, \cref{ch:takesaki}) $\N\rtimes_\theta\R\cong\M\bar\otimes\Bnd{L^2(\R)}$. Moreover, $\M$ is a factor of type III$_1$ iff $\N$ is a type II$_\infty$ \emph{factor}.
\end{theorem}
The slogan: \emph{every type III algebra sits inside a type II$_\infty$ algebra equipped with a trace-scaling flow, as its fixed points; the II$_\infty$ algebra is the type III algebra ``dressed by one clock''}, and type III$_1$ ones correspond to II$_\infty$ factors. In physical terms: add to the type III$_1$ algebra a variable that implements the modular flow (a ``clock'' $x$ and its conjugate, with $\Delta^{it}\leftrightarrow\ee^{-\ii tK}$), and a trace appears. This is the structure that will be realized physically by the observer's energy (\cref{ch:add_observer}) and, with gravity, will give the entropy (\cref{ch:clpw}). The counterpart of $\theta_s$ is shifts of the observer's energy; since the trace is $\ee^{-s}$-covariant under it, $\tau$ is infinite on the identity and the algebra is II$_\infty$, not II$_1$, until $q\ge0$ is imposed.

\begin{keyidea}
Type III arises from infinite systems in non-tracial states; the modular spectrum of ITPFI factors is the group generated by eigenvalue ratios. Connes' $S(\M)$ gives III$_0$, III$_\lambda$, III$_1$; local QFT algebras are III$_1$ (continuous modular spectrum). Connes--Takesaki: III$=$ II$_\infty\rtimes\R$, where $\R$ rescales the trace.
\end{keyidea}

\section*{Exercises}
\begin{exercise}
Compute the spectrum of $\Delta$ for $\bigotimes_j(M_2,\mathrm{diag}(p_j,1-p_j))$ with (a) $p_j=p$ fixed; (b) $p_j$ alternating between $p$ and $q$; (c) $p_j\to1/2$ summably, $\sum(p_j-\frac12)^2<\infty$. In which cases do you get type III$_\lambda$, III$_1$, and type II$_1$? (For (c), apply Kakutani, \cref{ch:inequiv}: the state is equivalent to the trace.)
\end{exercise}
\begin{exercise}
For the Powers factor with $\rho_j=\mathrm{diag}(p,1-p)$ and $\lambda=p/(1-p)$, show that $\Delta^{it_0}=\id$ on the GNS space for $t_0=2\pi/|\log\lambda|$, so $\sigma_{t_0}=\mathrm{id}$ is trivially inner. What about the flow of a different faithful state of the same algebra?
\end{exercise}
\begin{exercise}
For a free scalar field, the one-particle space of the right wedge is $L^2(\R,\dd\theta)$ in rapidity $\theta$, with boosts acting as translations $\theta\mapsto\theta+s$. Show the boost generator $B=-\ii\,\dd/\dd\theta$ has spectrum $\R$, so $\Delta=\ee^{-2\pi B}$ (on the one-particle space) has spectrum $[0,\infty)$. Compare with the discrete spectrum $\{\lambda^n\}$ of III$_\lambda$.
\end{exercise}
\begin{exercise}
Using only the existence of the trace $\tau$ on $\N$ with $\tau\circ\theta_s=\ee^{-s}\tau$, show $\tau(\id)=\infty$ for any $\N\ne0$. (If $\tau(\id)<\infty$ then $\tau(\theta_s(\id))=\tau(\id)=\ee^{-s}\tau(\id)$ forces $\tau(\id)=0$.) This is the sense in which the crossed product is II$_\infty$ and not II$_1$.
\end{exercise}
